\documentclass[letterpaper,12pt]{article}
\title{A Compact Measure for Working Papers}
\author{TeXPad / os8088}
\date{August 2026}
\begin{document}
\maketitle

\begin{abstract}
This note is a short working paper set in one monospaced face on
US Letter stock. The pad understands enough TeX for text, sections
and ruled tables. Colour is not used.
\end{abstract}

\section{Introduction}
A research note should be readable in hard copy. TeXPad keeps the
measure, the padding and the page number under the Layout menu, and
the same values may be written in the source.

Paragraphs wrap to the text width. A blank line starts a new
paragraph. \textbf{Bold} and \emph{emphasis} are the only faces.

\section{Method}
\subsection{Page}
Letter paper is 612 by 792 points. Normal margins are one inch on
every side. Narrow and wide steps are three-quarter and one-and-a-half
inches. Compact, normal and loose padding change paragraph skip,
leading and table cell inset.

\subsection{Tables}
Columns are \texttt{l}, \texttt{c} or \texttt{r}. A bar in the
preamble draws a rule. Cell padding follows the padding setting.

\begin{tabular}{|l|c|r|}
\hline
Part & Qty & Notes \\
\hline
8088 CPU & 1 & 4.77 MHz, 8-bit bus \\
8259A & 1 & PIC, master \\
2764 & 8 & 8K EPROM each \\
\hline
\end{tabular}

\begin{quote}
Quoted extracts are inset from both margins. They are meant for a
sentence taken from another page, not for decoration.
\end{quote}

\section{Results}
File, Export PDF and Export PS write a PDF 1.4 or a
PostScript Level 1 file, which a reader on another machine
opens. The body is Courier; rules are ordinary path strokes.

The book class adds chapter openings. Article is the default for a
note of this length.

\newpage
\section{Notes}
Special characters: 50\%, A \& B, \$32, file\_name. Line break\\
stays on the same paragraph.

\begin{itemize}
\item F5 typesets the preview.
\item Layout cycles class, margins and padding.
\item The system 8x8 font is the preview face.
\end{itemize}

\end{document}
